% =====================================================================
%  Disclosure: A Labour of Love
%  The philosophy of The Disclosure Platform. A stand-alone essay for a
%  general reader. Argues why an open, peer-reviewed, on-chain archive
%  democratizes the scientific method and builds an evidence-backed,
%  evolving consensus reality --- and why evidence is not proof. Written
%  to be read as an invitation: to read the archive, to file evidence, and
%  to earn your way into the named network that keeps it.
%  The engineering that operationalises this is described in a companion
%  paper, "The Evidence Peer-Review System".
%  Class : article
% =====================================================================

\documentclass[11pt,a4paper]{article}

% ---------- Packages ----------
\usepackage[T1]{fontenc}
\usepackage[utf8]{inputenc}
\usepackage{lmodern}
\usepackage[a4paper,margin=1in]{geometry}
\usepackage{microtype}
\usepackage{parskip}
\usepackage{titlesec}
\usepackage{enumitem}
\usepackage{booktabs}
\usepackage{array}
\usepackage{xcolor}
\usepackage{textcomp}
\usepackage[skins,breakable]{tcolorbox}
\usepackage{amsmath}
\usepackage{amssymb}
\usepackage{graphicx}
\usepackage[hidelinks,colorlinks=true,linkcolor=black,urlcolor=blue!55!black,citecolor=black]{hyperref}
\usepackage{fancyhdr}

% ---------- Colours ----------
\definecolor{ink}{HTML}{0E1116}
\definecolor{inkfaint}{HTML}{555B66}
\definecolor{accent}{HTML}{6A4CFF}
\definecolor{accent2}{HTML}{1F8F66}
\definecolor{warn}{HTML}{B85C00}
\definecolor{rule}{HTML}{C8CCD3}
\definecolor{panelbg}{HTML}{F6F4EE}

% ---------- Helper macros ----------
\newcommand{\arr}{\ensuremath{\rightarrow}}

% ---------- Page style ----------
\pagestyle{fancy}
\fancyhf{}
\fancyhead[L]{\footnotesize\textsc{The Disclosure Platform}}
\fancyhead[R]{\footnotesize\textsc{Disclosure --- A Labour of Love}}
\fancyfoot[C]{\footnotesize\thepage}
\renewcommand{\headrulewidth}{0.3pt}
\renewcommand{\footrulewidth}{0pt}

% ---------- Section style ----------
\titleformat{\section}
  {\normalfont\Large\bfseries\color{ink}}
  {\thesection}{0.6em}{}
\titleformat{\subsection}
  {\normalfont\large\bfseries\color{ink}}
  {\thesubsection}{0.6em}{}

% ---------- Helper boxes ----------
\newtcolorbox{plainbox}[1][]{
  enhanced, breakable, boxrule=0pt, frame hidden,
  colback=panelbg, colupper=ink,
  left=10pt, right=10pt, top=8pt, bottom=8pt,
  arc=2pt, #1
}

\newtcolorbox{notebox}[1][]{
  enhanced, breakable, boxrule=0pt, frame hidden,
  colback=panelbg, colupper=ink,
  left=12pt, right=12pt, top=8pt, bottom=8pt,
  borderline west={2pt}{0pt}{accent},
  arc=0pt, #1
}

% ---------- Title block ----------
\title{
  \vspace{-2em}
  {\Huge\bfseries Disclosure: A Labour of Love}\\[0.6em]
  {\Large The decentralized science (DeSci) network for contested\\
   truth --- an open archive that hands the scientific method\\
   back to the public and builds an evidence-backed, evolving reality}
}
\author{\textbf{The Disclosure Platform}}
\date{May 2026}

% =====================================================================
\begin{document}
% =====================================================================
\maketitle
\thispagestyle{empty}

\vspace{-1em}
\begin{plainbox}
\textbf{In one breath.}
The Disclosure Platform is a public record room for contested truth ---
\emph{decentralized science}, reclaimed for the thing science was supposed
to be \emph{for}: an open, evolving account of what we have reason to
believe, built and kept by the people who care about getting it right.
Anyone may read it. Anyone may file into it. But a claim never floats
free: it is filed \emph{against} the record that supports it --- a paper,
a book, a declassified file, a sworn testimony --- and a named network of
verified peers then judges that record in the open, by majority vote,
signing each verdict with their own cryptographic key, in a place no one
can quietly edit. No likes. No feed. No ads. No owner who can sell it,
seize it, or shut it down. This essay is the argument for why that matters
--- and an invitation to help build it. It is the scientific method handed
back to the public; it is honest about its own limits, because
\emph{evidence is not proof}; and it is, in the end, a labour of love.
\end{plainbox}

\tableofcontents
\vspace{1em}
\hrule
\vspace{1em}

% =====================================================================
\section{The truth needs a record room, not a feed}
% =====================================================================

We are not short of places to say things. We are drowning in them. What we
have lost is a place where what is said is \emph{kept} --- sourced, judged,
and answerable a year from now.

The modern feed is built to move. It rewards what is new, loud, and
agreeable; it forgets what is old, quiet, and inconvenient. A claim that
travels well is mistaken for a claim that is true, and the citation --- if
there ever was one --- falls away somewhere between the headline and the
share. When the conversation moves on, nothing remains that a later reader
could check. The feed is a river. You cannot step back into it, and it was
never built to let you.

A record room is the opposite of a feed. It is built to keep. Every entry
has a place, a source, and a history of who stood behind it and who stood
against it. Come back in a year and it is exactly where you left it, with
the reasoning still attached. The Disclosure Platform is a record room for
the kind of evidence the mainstream will not host --- not out of
contrarianism, but because a river was never the right instrument for
keeping a careful account.

\begin{notebox}
\textbf{The platform's one-line creed.} \emph{A feed forgets. A record
remembers.} Everything that follows is a consequence of taking that
sentence seriously.
\end{notebox}

This commitment shows up as a series of refusals, and they are not
sacrifices --- they are the whole point. There are no likes, because
popularity is not evidence. There is no ranking feed, because the order in
which the truth arrives is not the order of its importance. There is no
advertising and no data harvesting, because the moment a record room is
monetised by attention, it begins --- gently, then completely --- to
optimise for attention rather than for truth. There are no anonymous
moderators working out of sight, because a judgement no one can be named
for is a judgement no one can be answerable for. Strip all of that away and
what is left is \emph{decentralized science}: a public, peer-judged,
tamper-evident account of what we have reason to believe --- the part of
open inquiry that was always supposed to matter, finally given a home that
no single institution owns.

% =====================================================================
\section{The scientific method, democratized}
% =====================================================================

Science does not advance because scientists are unusually honest. It
advances because of a \emph{procedure} that does not depend on anyone's
honesty: make a claim, show your evidence, submit it to people who are free
to disagree, and let the claim stand only as long as it survives their
scrutiny. The genius of the method is not certainty. It is
self-correction. A mistake that enters the record can always be found and
removed by the same procedure that let it in.

The Disclosure Platform implements that exact loop, step for step --- and
throws every step open to anyone willing to do the work.

\begin{description}[leftmargin=1.2em,style=nextline]
  \item[\textnormal{\textbf{1. A claim, filed against a source.}}]
  In science a hypothesis means nothing without the data behind it. Here,
  every entry carries its citation. The claim cannot float; it has a record
  --- a paper, a book, a podcast, a declassified file, a deposition. The
  evidence is the entry. The assertion is just its label. And filing is
  open to everyone: you do not need permission to put a record on the table.

  \item[\textnormal{\textbf{2. Open peer review.}}]
  In science, work is judged by peers --- people competent and free to say
  no. Here, verified peers read the record and vote, in public, to canonize
  it, contest it, defend it, or remove it. The difference from the journal
  is that the room has no editor-in-chief and no closed door. The peers are
  named, their votes are signed, and the deliberation is visible \emph{while
  it happens}, not summarised after the fact.

  \item[\textnormal{\textbf{3. Challenge and replication.}}]
  In science, no result is ever truly closed; a finding can be re-examined
  the moment new evidence appears. Here, even an accepted record can be
  formally challenged, and the network must then decide again --- to keep it
  or to retire it. Nothing is canon forever by default; it is canon only for
  as long as it keeps surviving.

  \item[\textnormal{\textbf{4. Revision.}}]
  In science, the body of knowledge is whatever currently survives review.
  Here it is the same: the archive at any moment is the set of records the
  named network presently stands behind. When the evidence changes, the
  record changes with it.
\end{description}

What the platform democratizes is not the \emph{conclusion} --- you do not
get to vote a fact into being --- but the \emph{procedure}. The machinery of
careful judgement, historically locked inside institutions and journals, is
made into public infrastructure. The barrier to filing evidence is removed
entirely: anyone may submit. The responsibility of judging it is given to a
network that anyone may, over time, earn their way into. This is the
method's promise finally kept in the open: not ``trust us, we checked,''
but ``here is the evidence, here is who checked it, here is what they
signed, and here is how you can check them.''

\begin{notebox}
\textbf{Gatekeeping, replaced not abolished.} The platform does not pretend
that all judgement can be dispensed with --- a record room with no standards
is just a louder feed. It replaces \emph{closed} gatekeeping, where the gate
and the gatekeeper are hidden, with \emph{open} gatekeeping, where the gate
is a published rule and every keeper is named and accountable.
\end{notebox}

% =====================================================================
\section{Evidence is not proof}
% =====================================================================

Everything above could be misread as a machine for manufacturing
certainty. It is the opposite, and the difference is the moral centre of
the whole project. Get this part wrong and you have built another orthodoxy;
get it right and you have built something rarer than certainty --- something
honest.

\emph{Evidence is not proof.} A record supports a claim; it does not settle
it. A declassified memo is strong evidence that something was written down;
it is not proof that what was written is true. A sworn testimony is real
evidence that a person stated something under their name; it is not proof
that the thing stated occurred. To treat evidence as proof is to end the
inquiry exactly where it should continue.

The platform builds this humility into its structure rather than leaving it
to good intentions. Evidence is filed in tiers, and the tiers are honest
about their own strength:

\begin{itemize}[leftmargin=1.4em]
  \item \textbf{Tier I --- Peer-reviewed papers and declassified record.}
  The strongest footing. Documents that have already passed through some
  external scrutiny or official release.
  \item \textbf{Tier II --- Documented.} Books, institutional records, art
  --- evidence with a clear provenance but a lower bar of prior review.
  \item \textbf{Tier III --- Testimony.} First-person accounts, sworn or
  named. Indispensable, and the weakest as proof: a real record of what
  someone said, which is not the same as what is so.
\end{itemize}

A tier is not a verdict on truth. It is a statement about what \emph{kind}
of evidence you are looking at, so a reader can weigh it for themselves. The
platform never collapses that distinction. It does not display a green badge
that means ``true.'' What it can show is something more careful and more
honest: \emph{this record exists, here is its kind, and a majority of this
named network currently stands behind its place in the archive.}

\begin{notebox}
\textbf{What ``canon'' means here.} When a record is canonized it has not
been proven. It has been judged --- by a true majority of the named peers,
under a published, tier-dependent bar --- to belong in the archive on the
strength of its evidence, \emph{for now.} ``Canon'' means ``the network
presently stands behind this,'' never ``true forever.'' That single,
deliberate restraint is what separates a record of the truth from a monument
to an opinion.
\end{notebox}

This is also why the right to challenge can never be switched off. If
evidence were proof, a settled record could be sealed and forgotten. But
because evidence only ever \emph{supports}, any record --- however long it
has stood --- remains open to fresh evidence and renewed argument. The
inquiry is never declared finished, because honestly, it never is.

% =====================================================================
\section{An evidence-backed, evolving consensus reality}
% =====================================================================

If no single record is ever final, what does the platform actually build?
It builds a \emph{consensus reality}: not a list of certified facts, but a
living, structured, public account of what a named community has the best
current reason to believe --- with every reason attached, and every reason
revisable.

Two words in that phrase carry the weight.

\textbf{Evidence-backed.} Nothing enters the account on authority or
volume. Every node of it traces down to a citation a reader can follow. The
consensus is not a vote about how people feel; it is a judgement about what
the record will bear.

\textbf{Evolving.} The account is never finished, and it is not handed down.
It grows in two directions at once. It grows \emph{wider} as the community
proposes new \emph{pillars} --- whole domains of inquiry the network agrees
are worth keeping a record of. It grows \emph{deeper} as each pillar gains
\emph{topics}, the finer subjects beneath it under which specific evidence is
filed. And crucially, it can also \emph{shrink and correct}: a record that
no longer survives scrutiny can be retired, and one that is challenged and
re-defended is reaffirmed --- now stronger for having been tested.

There is no seeded ``official'' version of reality handed to the community
at the start. \textbf{The archive begins empty.} Every pillar, every topic,
every entry is brought into being by the peers themselves, through proposal
and consensus. The map is drawn by the people who walk the territory, and
redrawn whenever the territory turns out to be different from the map. The
blank page is not a flaw to apologise for --- it is the opportunity. What the
archive becomes is, quite literally, whatever its community is willing to
file and defend.

\begin{notebox}
\textbf{A map, not a monument.} A monument records what someone decided was
true and freezes it. A map records what we currently understand and expects
to be corrected. The Disclosure Platform is built to be a map: detailed
where the evidence is strong, honest about where it is thin, and always open
to redrawing. A consensus reality that cannot evolve is just an orthodoxy
with better production values.
\end{notebox}

This is what makes the result trustworthy in the only way a human account
of contested reality can be: not because it claims to be right, but because
it shows its work, names its reviewers, and keeps the door open for the next
correction. Trust is not asked for. It is made checkable.

% =====================================================================
\section{Named peers, and a record no one can capture}
% =====================================================================

A record room is only as honest as the people keeping it --- and only as
durable as its resistance to being seized. So the platform refuses three
conveniences that almost every other system accepts: anonymous authority, a
boss, and a kill switch.

\textbf{No anonymity for judges, and no unsigned acts.} Anyone may read
the archive, and anyone may file evidence into it --- those doors are
open to the world. But the people who \emph{judge} the evidence are
named, and \emph{every} act they take --- approving or rejecting a
record, challenging or defending it, proposing a new pillar or topic,
endorsing one, nominating a peer, voting on a removal, even removing a
captured owner --- is a cryptographic signature their wallet produces
and the contract itself verifies. There is no path by which the
platform, or anyone else, can post a vote on a peer's behalf, because
that call would not carry that peer's signature. There are no anonymous
moderators deleting records in the dark; there is no back office where
counts can be quietly adjusted. If a record is canonized, contested, or
retired, you can see who stood behind that decision and what they put
their name to. Accountability here is not a policy that can be quietly
suspended; it is the physical shape of the system.

\textbf{And every act now carries the reason for it.} Alongside the
signature, each peer may attach a short deliberation note --- the
\emph{why} behind their yes, no, defend, challenge, or endorse. The note
is bound into the signature itself (its fingerprint is part of what the
peer signs), so a peer cannot quietly walk back the reasoning they were
counted for: the record now keeps not only \emph{who} voted and
\emph{how}, but \emph{why}, in their own words, under their own name.
``Judgement under your own name'' has become ``judgement under your own
signature, with your reasons attached, for anyone to read.''

\textbf{No administrator in the loop.} Membership in the named network is
not granted by an owner. A newcomer is nominated by an existing peer and
must then earn the endorsements of others before the system itself --- by
its own published rule --- recognises them. No one can appoint a loyalist; no
one can expel a dissenter by decree. The community admits and, if it must,
removes its own members by vote.

\textbf{No one can capture it --- not even its author.} The platform launches
from a small founding set of peers, and in that earliest phase a caretaker
role exists only to seed the first members and to pause in a genuine
emergency. It can never force a record in, force a peer out, or override a
vote. And once the network can stand on its own, that role can be
\emph{renounced} entirely, leaving no privileged account at all. Should a
caretaker ever turn --- freeze the system and refuse to let go --- the peers
themselves can strip the role by a two-thirds supermajority and carry on.
The keys do not pass to a new owner; they pass to no one. This is what
``owned by no one'' means when you build it instead of merely promising it.

The effect is a community of people who carry handles and signatures, not
followers and vanity metrics. There is no profile photo to curate, no
follower count to chase, no engagement loop to feed. A peer's standing comes
from the records they have helped verify and the judgements they have been
willing to sign --- the only reputation a record room has any business
keeping.

% =====================================================================
\section{The destination: Disclosure, on every front}
% =====================================================================

It is worth saying out loud what all of this is \emph{for}. The name on
the door is not decorative. The Disclosure Platform is built to become
the place where \textbf{Disclosure} --- on every front where institutions
have for decades decided what the public is allowed to know --- finally
has somewhere to live, to accumulate, and to stand.

\emph{Disclosure} here is plural by design. Whatever the file room of a
ministry, an agency, a corporation, a university, a publisher, or a
priesthood has kept off the public record, this archive is the place to
put it. A declassified intelligence document. A whistleblower's sworn
testimony. A buried preprint. A suppressed clinical dataset. A
photographed artefact. A historian's primary source the textbook left
out. A first-person account a court would not hear. Wherever a
gatekeeper has been the one deciding what counts, an open record room
is the answer to the gatekeeper.

What is unusual about this platform is not the ambition --- many places
have aspired to host the unhostable --- but that the architecture is
finally adequate to the ambition. Every previous attempt to keep a
public record of contested truth has eventually hit the same three
walls: the host gets pressured, the moderators get captured, or the
funding gets pulled. Each of those failure modes has a specific
engineered answer here, and the answers do not merely promise
resistance, they \emph{enforce} it.

\begin{itemize}[leftmargin=1.4em]
  \item \textbf{The host cannot be pressured into pulling a record},
  because there is no host with that power. The canonical record lives
  on a public blockchain; the only way to remove it is the same way it
  entered: a vote, in the open, by a supermajority of named peers,
  under their own signatures. ``Take it down'' has no recipient.

  \item \textbf{The moderators cannot be captured}, because there are
  no moderators in the closed-door sense. Every act of admission,
  removal, ratification, retirement, canonization, expulsion, or
  challenge is a cryptographic signature from a named wallet that the
  contract itself verifies. A coalition of less than two thirds of
  peers can do nothing the rest do not consent to; a coalition of two
  thirds can even strip a captured caretaker entirely. The most
  consequential power in the system --- evicting its own owner --- has
  been engineered to be a signed, counted, public act, not a
  back-channel deal.

  \item \textbf{The funding cannot be pulled}, because there is no
  funding to pull. The platform sells nothing, harvests nothing, and is
  owned by no one who could be paid to fold it. Every state change is
  paid for by the wallet that causes it --- the user, never the
  platform --- so the operating cost of the archive is borne by the
  people using it, and there is no central treasury an adversary could
  drain or seize.
\end{itemize}

This is what makes ``Disclosure on every front'' not a slogan but a
plausible outcome. The instrument is correctly shaped for the job. An
archive that cannot be edited, cannot be captured, and cannot be
defunded is the right kind of place for evidence that powerful
interests would prefer to see disappear. It is also the right kind of
place for evidence that powerful interests have, until now, simply been
\emph{able} to make disappear, because there was nowhere else for it to
go.

\begin{notebox}
\textbf{The blank page is an invitation, not a verdict.} The archive
starts empty on purpose. The platform does not arrive with a settled
view of what has been hidden and where; it arrives with the
\emph{capacity} to keep an honest, signed, evolving record of whatever
the network is willing to file and defend. Which fronts of disclosure
become real here is, in the most literal sense, up to the peers who
show up to keep them. Every domain where the truth has been deferred is
a pillar waiting to be proposed.
\end{notebox}

None of this means the platform decides what is true; it never has and
it never will. ``Canon'' here still means only ``a majority of named
peers presently stands behind this record, on the strength of its
evidence.'' Disclosure, in this archive, is not the end of inquiry ---
it is the moment a piece of evidence finally has somewhere to be filed,
judged in the open, and kept where it cannot quietly disappear again.
That is a smaller claim than ``the truth has been revealed,'' and it is
the only claim a record room of this kind can honestly make. It is also
exactly the claim that, scaled across enough domains and enough years,
\emph{is} what disclosure looks like in practice: not one dramatic
revelation, but a thousand records the public was never supposed to
have, kept somewhere they cannot be erased, with the names of the
people who stood behind them attached.

% =====================================================================
\section{What this is not}
% =====================================================================

It is worth being plain about the things the platform deliberately does not
do, because each absence is a design decision in service of the record ---
and, taken together, a quiet manifesto.

\begin{center}
\begin{tabular}{@{}ll@{}}
\toprule
\textbf{Not present} & \textbf{Because} \\
\midrule
No likes            & Popularity is not evidence. \\
No feed             & A river forgets; a record keeps. \\
No ranking          & Importance is not a function of recency. \\
No algorithm        & Nothing decides for you what is true. \\
No followers        & Standing comes from signed judgement, not reach. \\
No ads              & Attention must not become the product. \\
No data collection  & Readers are not the inventory. \\
No anonymous mods   & Judgement must be answerable. \\
No revenue model    & The record is not for sale. \\
No owner            & The keys belong to no one. \\
No closed source    & The gate must be inspectable by all. \\
\bottomrule
\end{tabular}
\end{center}

\vspace{0.5em}

The Disclosure Platform is infrastructure for evidence, not a place to
socialise in the old, hollow sense. Submissions are public the instant they
land. The contract that holds the rules is the only gatekeeper, and that
contract is open. Anyone may read; named peers verify; the public watches.
That is the whole arrangement, and its plainness is the point.

% =====================================================================
\section{Your part in this}
% =====================================================================

A record room is not a spectator sport. It is built --- and kept honest ---
by the people who show up. There is a place in it for you whatever you bring,
and the door is genuinely open at every level.

\begin{description}[leftmargin=1.2em,style=nextline]
  \item[\textnormal{\textbf{Read it.}}]
  The archive is public, free, and owned by no one. Follow any claim down to
  its source. Re-check any peer's signature in your own browser and confirm
  that the named person --- not the platform --- authored that vote. Trust
  nothing on our say-so; verify it yourself. That option, always available,
  is the product.

  \item[\textnormal{\textbf{File evidence.}}]
  Filing is open to everyone --- no application, no gatekeeper. Found a paper,
  a document, a testimony the mainstream won't keep? Put it on the table,
  filed against its source and its tier, and let the named network judge it
  in the open. The archive only grows because people bring it something to
  weigh.

  \item[\textnormal{\textbf{Become a peer.}}]
  The network of judges is not a closed guild --- it is something you earn
  your way into. Get nominated by an existing peer, win the endorsements of
  others, and the system itself recognises you. No résumé, no owner's
  blessing. From there you canonize evidence, propose new pillars and topics,
  and defend or challenge the record under your own name.

  \item[\textnormal{\textbf{Keep it honest.}}]
  Canon is never sealed. If a standing record no longer deserves to stand,
  challenge it. If a good record is unfairly attacked, defend it. The
  archive's integrity is not maintained by an institution; it is maintained
  by peers willing to keep doing the work.
\end{description}

\begin{notebox}
\textbf{Why now.} The archive is at its blank-page moment --- the rarest and
most consequential time to arrive. The first pillars, the first topics, the
first canonized records, the founding shape of the named network: all of it
is still unwritten. Early hands do not just add to the record; they set its
standards. There has rarely been a better time to help decide what a public
record of contested truth should look like.
\end{notebox}

% =====================================================================
\section{A labour of love}
% =====================================================================

Strip away the things the platform refuses --- the advertising, the data
harvesting, the engagement metrics, the paywall, the owner --- and an obvious
question remains: why would anyone build this, and why would anyone keep it?

There is no financial answer. No one is enriched by a canonized record. The
peers are not paid to read evidence or to sign their judgements; they give
their attention and their names to keep the account honest. The archive
sells nothing and harvests no one. By every metric the modern internet is
built to maximise, the platform is a failure on purpose --- and that failure
is exactly the freedom that lets it serve the truth instead of the market.

What is left, once profit is removed as a motive, is the oldest reason to
keep a careful record at all: because the truth deserves a room of its own,
and because building that room for other people --- most of whom you will
never meet --- is an act of care.

That is what we mean by a labour of love. Not sentiment, but the specific
kind of love that shows up as patient, unpaid, accountable work: the
willingness to write things down properly, to source them, to argue about
them in the open under your own name, and to let your own conclusions be
overturned when the evidence turns. It is the love of getting it right over
the comfort of being right. It is the conviction that an honest, open,
evolving record of what we have reason to believe is worth making for its
own sake --- and worth giving away.

\begin{notebox}
\textbf{The whole thing, finally.} A public record room for contested
truth, where evidence is filed against its source, judged in the open
by a majority of named peers --- each verdict cryptographically signed
with the reasoning attached --- kept where no one can quietly edit it,
owned by no one and seizable by no one, and revised forever as the
evidence changes. The scientific method, handed back to the public.
Decentralized science --- a DeSci network reclaimed for the truth. The
instrument, finally adequate to the ambition: Disclosure, on every
front, with the evidence to back it and the signatures to stand behind
it. A labour of love --- and an open invitation to help build it.
\end{notebox}

The companion paper, \emph{The Evidence Peer-Review System}, describes the
engineering that makes every promise in this essay enforceable rather than
merely stated --- how ``filed against a source'' becomes a tamper-evident
hash, how ``judged by named peers'' becomes a signed and counted vote, how
``owned by no one'' becomes a contract a rogue hand cannot hold, and how
``revised forever'' becomes a lifecycle no single actor can override. The
mechanism is the proof that the philosophy is meant.

\end{document}
